\section{Occupation and remaining duration}
\label{sec:occupation}

The next estimate is the robust analogue of the critical occupation bound
in \citet{OpenAI2026ThreeSATVariance}.  It uses the lower conditional-density
bound for each ordinary clause.

\begin{lemma}[Conditional occupation estimate]
\label{lem:occupation}
There is a constant $K_0=K_0(B_{\mathrm{o}},k,\rho)<\infty$ such that, for every
$i\le m<L$,
\begin{equation}
    \E\left[
      \sum_{s=m}^{L-1}
        \ind_{\{S_s\ne\varnothing\}}p_s^\eta
      \,\middle|\,\cH_m
    \right]
    \le K_0d^\eta.
    \label{eq:occupation-bound}
\end{equation}
\end{lemma}

\begin{proof}
Suppose $s+1$ is ordinary, a fact known under $\cE$.  Then
$S_{s+1}=S_s[C_{s+1}]$, and the hidden contents of $C_{s+1}$ remain fresh
given $\cH_s$.  By the ordinary-clause density hypothesis, their conditional law is
bounded below by $B_{\mathrm{o}}^{-1}U_{u,k}$.  Apply
Lemma~\ref{lem:robust-potential} conditionally to obtain
\begin{align}
    \ind_{\{S_s\ne\varnothing\}}p_s^\eta
    \le{}&
    kB_{\mathrm{o}}\,d^{\eta-1}
      \E\bigl[P_d(S_{s+1})-P_d(S_s)\mid\cH_s\bigr]
      \notag\\
    &+c_kd^\eta u^{-k}.
    \label{eq:conditional-drift}
\end{align}
Condition on $\cH_m$ and sum over ordinary steps.  Every potential
increment is nonnegative, including increments at nonordinary steps.  Thus,
pathwise,
\[
    \sum_{\substack{m\le s<L\\s+1\text{ ordinary}}}
       \bigl(P_d(S_{s+1})-P_d(S_s)\bigr)
    \le P_d(S_L)-P_d(S_m)
    \le d.
\]
The ordinary-step contribution to \eqref{eq:occupation-bound} is therefore
at most
\[
    kB_{\mathrm{o}}\,d^\eta+c_kLd^\eta u^{-k}.
\]
There are at most $Z$ remaining steps, because every nonordinary index
strictly after $i$ meets at least one root.  Each remaining summand is at
most one.  Since $Z\le d\le d^\eta$, $L\le\rho n$, and $u=n-k\ge n/2$,
we obtain
\[
    \E\left[
      \sum_{s=m}^{L-1}
        \ind_{\{S_s\ne\varnothing\}}p_s^\eta
      \,\middle|\,\cH_m
    \right]
    \le
    \left(kB_{\mathrm{o}}+1+c_k\rho 2^k\right)d^\eta.
\]
This proves the assertion with, for example,
$K_0=kB_{\mathrm{o}}+1+c_k\rho2^k$.
\end{proof}

\begin{corollary}[Remaining duration]
\label{cor:remaining-duration}
For $i\le m<L$,
\begin{equation}
    \E\left[
      \sum_{s=m}^{L-1}\ind_{A_s}
      \,\middle|\,\cH_m
    \right]
    \le W_m,
    \label{eq:remaining-duration}
\end{equation}
where
\begin{equation}
    W_m=
    \begin{cases}
      \min\{L,K_0d^\eta p_m^{-\eta}\},&p_m>0,\\
      L,&p_m=0.
    \end{cases}
    \label{eq:W-definition}
\end{equation}
In particular, $W_m$ is $\cG_m$-measurable.
\end{corollary}

\begin{proof}
For $s\ge m$, one has $p_s\ge p_m$, and $A_s$ implies
$S_s\ne\varnothing$.  If $p_m>0$, then pathwise
\[
    \sum_{s=m}^{L-1}\ind_{A_s}
    \le
    p_m^{-\eta}
    \sum_{s=m}^{L-1}
      \ind_{\{S_s\ne\varnothing\}}p_s^\eta.
\]
Apply Lemma~\ref{lem:occupation}.  The count is always at most $L$, which
also gives the case $p_m=0$.  Measurability follows from that of $d$ and
$p_m$.
\end{proof}
